0382: \phantomsection\label{struct:V4-C06-CH14}
0383: % KNOWLEDGE-PLAN: KN-V4-C29-CONCEPT-009 L1
0384: \paragraph{读前自检：几何解释：每个 k 是 (M-1) 维单词单纯形中的主题点。}PLSA数值例给出$\theta_{:j}=(0.6,0.4)^{\mathsf T}$及词概率$(0.2,0.7)$；请算$P(w_i\mid d_j)=0.4$和两主题后验，核对责任度为$(0.3,0.7)$且和为1。\par
0385: 
0386: \begin{geometrybox}
0387: \textbf{几何解释。}每个$\phi_{:k}$是$(M-1)$维单词单纯形中的主题点。文档分布$P(w\mid d_j)=\Phi\theta_{:j}$是这些主题点的凸组合，因此全部训练文档都落在由主题点张成的凸包中。$K=1$时凸包退化为一点；$K=2$时为线段；主题点线性相关或重复时，有效维数低于$K-1$。
0388: \end{geometrybox}
0389: 
0390: \phantomsection\label{struct:V4-C06-CH15}
0391: \begin{probabilitybox}
0392: 责任度$\gamma_{ijk}$不是“单词$w_i$永远属于主题$k$”的概率，而是给定单词与文档后的局部后验。同一个词在不同文档中可由不同主题解释。$\Phi$刻画主题的词分布，$\Theta$刻画训练文档的主题分布，二者必须与后验责任度分开。
0393: \end{probabilitybox}
0394: 
0395: \phantomsection\label{struct:V4-C06-CH16}
0396: \begin{optimizationbox}
0397: \textbf{优化解释。}固定经验$\pi_j$时，PLSA等价于在列随机约束下最小化经验条件分布与$\Phi\Theta$之间的加权KL散度。目标对$\Phi$和$\Theta$分别具有较好结构，但对二者联合非凸；主题置换、重复主题和边界零概率都会造成非唯一解。
0398: \end{optimizationbox}
0399: 
0400: \phantomsection\label{struct:V4-C06-CH17}
0401: % FIGURE-KNOWLEDGE: fig:V4-C06-simplex=SLM-18-01-005
0402: % v2.6.0 figure UID: FIG-P525-01
\tikzset{
  slfig-FIG-P525-01/.style={font=\fontsize{9.4pt}{11.2pt}\selectfont,
    every path/.append style={line cap=round,line join=round}},
  slfig-FIG-P525-01-outer/.style={draw=SLTextGray,line width=.92pt},
  slfig-FIG-P525-01-hull/.style={draw=SLTeal,line width=.86pt},
  slfig-FIG-P525-01-spoke/.style={draw=SLGray,densely dashed,line width=.54pt},
  slfig-FIG-P525-01-weight/.style={fill=none,inner sep=0pt},
  slfig-FIG-P525-01-doclabel/.style={anchor=west,text=SLGold,fill=white,
    fill opacity=.92,text opacity=1,inner xsep=1pt,inner ysep=.5pt},
  slfig-FIG-P525-01-formula/.style={draw=SLRuleGray,rounded corners=2pt,
    fill=SLSoftGray,line width=.62pt,align=center,inner ysep=2.2pt},
  slfig-FIG-P525-01-legend/.style={anchor=west,text=SLTextGray,
    font=\fontsize{8.8pt}{10.4pt}\selectfont}
}
\begin{figure}[H]
\centering
\begin{tikzpicture}[slfig-FIG-P525-01,x=1cm,y=1cm]
  % 显式重心变换：(p_1,p_2,p_3) ->
  % (x,y)=(6(p_2+p_3/2), 6(sqrt(3)/2)p_3)。
  \newcommand{\barypoint}[4]{%
    \coordinate (#1) at ({6*(#3+0.5*#4)},{6*0.8660254*#4});}
  \coordinate (Wone) at (0,0);
  \coordinate (Wtwo) at (6,0);
  \coordinate (Wthree) at (3,{3*sqrt(3)});
  \barypoint{T1}{.78}{.12}{.10}
  \barypoint{T2}{.10}{.78}{.12}
  \barypoint{T3}{.12}{.13}{.75}
  \barypoint{Doc}{.309}{.4195}{.2715}

  \fill[SLTeal,opacity=.085] (T1)--(T2)--(T3)--cycle;
  \draw[slfig-FIG-P525-01-outer] (Wone)--(Wtwo)--(Wthree)--cycle;
  \node[below left] at (Wone) {$w_1$};
  \node[below right] at (Wtwo) {$w_2$};
  \node[above] at (Wthree) {$w_3$};
  \draw[slfig-FIG-P525-01-hull] (T1)--(T2)--(T3)--cycle;

  \foreach \p/\lab/\pos in {T1/$\phi_{:1}$/left,T2/$\phi_{:2}$/right,T3/$\phi_{:3}$/above right}{
    \fill[SLBlue] (\p) circle (2.6pt) node[\pos,text=SLBlue] {\lab};
  }
  \draw[slfig-FIG-P525-01-spoke] (Doc)--(T1);
  \draw[slfig-FIG-P525-01-spoke] (Doc)--(T2);
  \draw[slfig-FIG-P525-01-spoke] (Doc)--(T3);
  \node[slfig-FIG-P525-01-weight] at (2.05,1.28) {$\theta_{1j}$};
  \node[slfig-FIG-P525-01-weight] at (3.65,.95) {$\theta_{2j}$};
  \node[slfig-FIG-P525-01-weight] at (2.55,2.05) {$\theta_{3j}$};
  \node[diamond,draw=SLGold,fill=SLGold!18,minimum size=6mm,
    inner sep=0pt,line width=.9pt] (dj) at (Doc) {};
  \node[slfig-FIG-P525-01-formula,
    minimum width=58mm,minimum height=8mm] at (8.0,2.55)
    {$P(w\mid d_j)=\sum_{k=1}^{3}\theta_{kj}\phi_{:k}$\\
     $\theta_{kj}\ge0,\ \sum_k\theta_{kj}=1$};
  \node[slfig-FIG-P525-01-legend] at (6.95,1.25)
    {圆：主题点\quad 菱形：文档分布 $P(w\mid d_j)$};
\end{tikzpicture}
\caption{三个单词维度中的主题单纯形：文档分布是主题点按$\theta_{:j}$形成的凸组合}\label{fig:V4-C06-simplex}
\end{figure}

0403: 图\ref{fig:V4-C06-simplex}把外层三角形表示为单词单纯形，把蓝色三角形表示为主题凸包。文档点的位置由$\theta_{:j}$决定；不同参数可以给出相同文档点，说明潜在表示不必唯一。
0404: 
0405: \phantomsection\label{struct:V4-C06-CH18}
0406: \textbf{小型数值例子。}设某文档只有两个主题，$\theta_{:j}=(0.6,0.4)^{\mathsf T}$；某单词在两主题下的概率为$0.2$与$0.7$。则
0407: $P(w_i\mid d_j)=0.2\times0.6+0.7\times0.4=0.4$。若观察到该词，后验责任度为
0408: \[
0409: P(z_1\mid w_i,d_j)=\frac{0.2\times0.6}{0.4}=0.3,
0410: \qquad
0411: P(z_2\mid w_i,d_j)=0.7.
0412: \]
0413: 先验主题权重为$(0.6,0.4)$，观察到更偏向第二主题的词后，后验转为$(0.3,0.7)$。
